\documentclass[12pt, french, latin.medieval, a4paper]{report}
% package de langues 
\usepackage[french]{babel}
\usepackage[utf8]{inputenc}
\usepackage[T1]{fontenc}

% package pour référence de liens sur pdf 
\usepackage[pdftex]{hyperref} % liens
\usepackage{lastpage} % numéro de page
\usepackage{caption}

% package de formes de document 
\usepackage[left=2.5cm,
              right=2.5cm,
              top=2.5cm,
              bottom=2.5cm
             ]{geometry}
% pour les tableaux 
\usepackage{array,multirow,makecell} 
\newcolumntype{R}[1]{>{\raggedleft\arraybackslash }b{#1}}
\newcolumntype{L}[1]{>{\raggedright\arraybackslash }b{#1}}
\newcolumntype{C}[1]{>{\centering\arraybackslash }b{#1}}

% package pour la page de couverture 
\usepackage{graphicx}
\usepackage[cyr]{aeguill}
\usepackage{fancyheadings}
\usepackage{amsmath}
\usepackage{amsfonts}
\usepackage{amssymb}
\usepackage{geometry}
\usepackage{eso-pic}
\usepackage{tikz}
\usetikzlibrary{calc}


% nouvelle commande pour la forme 
%\fancyhf[OFC]{\thepage / \pageref{LastPage}}
\rfoot{\thepage / \pageref{LastPage}}
%\setcounter{tocdepth}{3}

% pakages de couleurs du document
\usepackage{color}
\usepackage{xcolor}


% package pour la chimie
\usepackage[version=4]{mhchem}


% package pour la visualisation du code
\usepackage{listings}
\definecolor{darkWhite}{rgb}{0.94,0.94,0.94}
 
\lstset{
  aboveskip=3mm,
  belowskip=-2mm,
  backgroundcolor=\color{darkWhite},
  basicstyle=\footnotesize,
  breakatwhitespace=false,
  breaklines=true,
  captionpos=b,
  commentstyle=\color{red},
  deletekeywords={...},
  escapeinside={\%*}{*)},
  extendedchars=true,
  framexleftmargin=16pt,
  framextopmargin=3pt,
  framexbottommargin=6pt,
  frame=tb,
  keepspaces=true,
  keywordstyle=\color{blue},
  language=python,
  literate=
  {²}{{\textsuperscript{2}}}1
  {⁴}{{\textsuperscript{4}}}1
  {⁶}{{\textsuperscript{6}}}1
  {⁸}{{\textsuperscript{8}}}1
  {€}{{\euro{}}}1
  {é}{{\'e}}1
  {è}{{\`{e}}}1
  {ê}{{\^{e}}}1
  {ë}{{\¨{e}}}1
  {É}{{\'{E}}}1
  {Ê}{{\^{E}}}1
  {û}{{\^{u}}}1
  {ù}{{\`{u}}}1
  {â}{{\^{a}}}1
  {à}{{\`{a}}}1
  {á}{{\'{a}}}1
  {ã}{{\~{a}}}1
  {Á}{{\'{A}}}1
  {Â}{{\^{A}}}1
  {Ã}{{\~{A}}}1
  {ç}{{\c{c}}}1
  {Ç}{{\c{C}}}1
  {õ}{{\~{o}}}1
  {ó}{{\'{o}}}1
  {ô}{{\^{o}}}1
  {Õ}{{\~{O}}}1
  {Ó}{{\'{O}}}1
  {Ô}{{\^{O}}}1
  {î}{{\^{i}}}1
  {Î}{{\^{I}}}1
  {í}{{\'{i}}}1
  {Í}{{\~{Í}}}1,
  morekeywords={*,...},
  numbers=left,
  numbersep=10pt,
  numberstyle=\tiny\color{black},
  rulecolor=\color{black},
  showspaces=false,
  showstringspaces=false,
  showtabs=false,
  stepnumber=1,
  stringstyle=\color{gray},
  tabsize=4,
  title=\lstname,
}

%package pour les algorithmes
\usepackage{algorithm}
\usepackage{algorithmic}
\usepackage{amsmath}

% package pour les tableaux
\usepackage{longtable}
\usepackage{booktabs}

% package pour les figures 
\usepackage{subcaption}

% package de reférencement
% \usepackage[backend=biber,safeinputenc, sorting=none]{biblatex}

% mettre la bibliographie dans la table de matières
%\usepackage[nottoc, notlof, notlot]{tocbibind}
\include{page_couverture.tex}


\author{MBE MBE MINDJANA LOIC HENRI}
\title{GNPS et MDByTE}

\begin{document}
    
    \pagecouverture
    {Outils bio-informatiques GNPS et MADByTE}
    {}
    \tableofcontents
    \vspace{20cm}

    \section{Informations sur la séance}
    \begin{itemize}
        \item Date : \textbf{20 Avril 2023}, \textbf{13h}(date locale)
        \item Moyen : visioconférence
    \end{itemize}
    
    \section{Sujet de la séance}
    Il était question pendant cette séance que nous présentions les études faites (qui ont été sujet de la rédaction d'un rapport) sur les outils GNPS et MADByTE sur leurs fonctions et leurs fonctionnements, pour avoir une meilleure orientation sur le travail.
    
    \section{Taches réalisées lors de la séance}
    Nous avons présenter l'outil MADByTe qui manipule les données TOCSY et HSQC de molécules et GNPS qui manipule les données $MS^{2}$ de molécules pour réaliser la déréplication de mélanges complexes suivant notre compréhension; et ainsi nous avons parler des procéssus qu'utilisent MADByTE et GNPS pour le faire. 
    
    \section{Taches pour la séance prochaine}
    A la fin de la séance Prof Waffo, nous a signaler qu'une reunion par visioconférence sera faite pour nous présenter les données $MS^{2}$, HSQC et TOCSY d'un point de vue d'un biologiste, dans le but qu'on les comprenne mieux, pour éfficassement combiner ces données pour réaliser de la déréplication.

\end{document}